\chapter{Stage 7: The Output Contract}
\label{ch:output}

Six stages have decided what to answer and how much to trust it. Stage 7 is the last
one in the per-query pipeline, and it is the only one the user ever sees. Its job is
narrow but important: take the four-outcome decision of \Cref{ch:composite} and turn
it into what is actually shown, and crucially, ensure the answer never travels alone.
This is the stage where the calibration mismatch of \Cref{ch:problem}, the confident
lie wearing the same clothes as the confident truth, is finally repaired at the
surface the user touches.

\begin{intuition}[The decision becomes a presentation]
By the time a query reaches Stage 7, the hard thinking is done: there is a generated
string and a verdict (store, defer, abstain, or discard). Stage 7 does not re-judge
anything. It decides what string to display, whether to show it at all, and what
reliability signal to attach. The guiding principle is that a user should never
receive a fluent answer with no indication of how much to trust it; the answer and
its confidence are delivered together, as one object.
\end{intuition}

\section{Two answers, kept separate}
\label{sec:output-two-answers}

The pipeline carries the answer in \emph{two} fields, and the separation is
deliberate. One field, the \emph{raw answer}, is the model's literal output, kept
verbatim. The other, the \emph{display answer}, is the curated string a user would
see. They are computed from the same generation but serve opposite masters.

\begin{intuition}[Why scoring and display must not share a string]
Evaluation needs the model's \emph{literal} prediction: to score exact match
honestly, the harness must see exactly what the model produced, including a wrong
answer, so that a wrong answer counts as wrong. The user, by contrast, should see a
\emph{curated} string: the final answer with its reasoning scaffold stripped away, or
an honest refusal, or nothing at all. If display and scoring shared one field, then
either the user would see raw scaffold text, or the scorer would be fed a cleaned-up
string that no longer matches what the model actually said. Keeping the raw answer for
the scorer and the display answer for the user lets each be exactly what it should be.
\end{intuition}

\section{Mapping the decision to a display string}
\label{sec:output-map}

The display answer is computed by one small function that reads the verdict and
nothing else. Its logic is the whole contract.

\begin{lstlisting}[caption={The display mapping, condensed from \texttt{caem/pipeline.py}.}, label={lst:display}]
if verifier_output is None:                # Tier 1 hit, or a pipeline edge case
    return strip_scaffold(answer)          # serve the (stored) answer as-is
if decision == "ABSTAIN":
    return "I do not know."                # the calibrated refusal
if decision == "DISCARD":
    return ""                              # suppress: the prediction is not safe to show
return strip_scaffold(answer)              # STORE / DEFERRED: show the final answer
\end{lstlisting}

So the four verdicts map to three surface behaviours. A \textbf{stored} or
\textbf{deferred} answer is shown, because the system trusts it enough to serve (and,
for a deferred answer, also holds it for reconsideration). An \textbf{abstained} query
returns the fixed refusal string ``I do not know.'', the principled failure mode in a
factual setting where a confident wrong answer costs more than an honest refusal. A
\textbf{discarded} answer is suppressed entirely, because the verdict was that the
prediction is not safe to expose. A Tier 1 memory hit, which skipped the verifier by
design, serves its stored answer directly.

\begin{precise}[Scaffold stripping keeps display and scoring aligned]
The generated string is the scaffolded form ``Reasoning: \dots\ Answer: $Y$'' from
\Cref{ch:tiers}. The display function strips this down to just the final answer $Y$,
and it does so through the \emph{same} extraction routine the exact-match scorer uses.
This shared routine is not an accident: it guarantees that the answer the user sees
and the answer the scorer grades are extracted identically, so a reported score always
corresponds to what was actually shown. If the string carries no scaffold markers (a
direct answer, or a legacy format), it is passed through unchanged.
\end{precise}

\section{The answer never travels alone: the confidence envelope}
\label{sec:output-envelope}

A bare string, even the right one, would reproduce the original problem: the user
could not tell a trusted answer from a guess. So a deployment frontend wraps the
display answer in a \emph{confidence envelope} built from the same verdict and the
calibrated score $\ustored$. The envelope carries six things alongside the answer.

\begin{description}[leftmargin=0pt, labelsep=0.5em, style=nextline]
  \item[A continuous score.] The calibrated $\ustored$ rendered as a percentage, for
  example ``72\%'', read directly as the chance the answer is correct.
  \item[A semantic tag.] One of \emph{verified}, \emph{provisional},
  \emph{conflicting}, or \emph{insufficient}, naming \emph{which} situation produced
  the verdict rather than just how confident it was.
  \item[An ordinal label.] One of \emph{high}, \emph{moderate}, \emph{low}, or
  \emph{very low}, an at-a-glance reliability hint.
  \item[An icon.] A check, a tilde, a warning, or a cross, for instant visual reading.
  \item[A show-or-hide flag.] Whether to display the answer at all, hidden on an
  abstention.
  \item[A caveat.] An optional hedge sentence appropriate to the verdict.
\end{description}

\begin{table}[htbp]
\centering
\caption[The confidence envelope]{How each verdict is rendered. The four verdicts map
to four distinct user-facing presentations, not to four points on a single confidence
scale.}
\label{tab:envelope}
\small
\renewcommand{\arraystretch}{1.25}
\begin{tabularx}{\textwidth}{@{}l l l l >{\raggedright\arraybackslash}X@{}}
\toprule
\textbf{Verdict} & \textbf{Tag} & \textbf{Label} & \textbf{Icon} & \textbf{Shown to the user} \\
\midrule
STORE & verified & high & check & the answer, with its confidence percentage \\
DEFERRED & provisional & moderate & tilde & the answer, with a ``re-checking later'' caveat \\
DISCARD & conflicting & low & warning & a hedge that the evidence found disagrees \\
ABSTAIN & insufficient & very low & cross & ``I do not know.'' (the answer is hidden) \\
\bottomrule
\end{tabularx}
\end{table}

\begin{figure}[htbp]
\centering
\begin{tikzpicture}[
  font=\small\sffamily,
  card/.style={rectangle, rounded corners=5pt, line width=1.2pt, align=left,
    inner sep=7pt, text width=6.2cm, minimum height=2.3cm,
    drop shadow={shadow xshift=0.4ex, shadow yshift=-0.4ex, fill=black!12}},
]
  % STORE / verified
  \node[card, draw=cbExample!70!black, fill=cbExample!6] (c1) at (0,0) {
    {\large\color{cbExample!55!black}$\checkmark$}\;\textbf{\color{cbExample!50!black}verified}
    \hfill {\color{cbInk!60}\footnotesize high $\cdot$ 70\%}\par\vspace{5pt}
    {\footnotesize\color{cbInk!70}\itshape Roquefort Cheese is made from what sort of milk?\par}\vspace{4pt}
    {\large sheep's milk.\par}};
  % DEFERRED / provisional
  \node[card, draw=cbScenario!70!black, fill=cbScenario!6] (c2) at (7.4,0) {
    {\large\color{cbScenario!55!black}$\sim$}\;\textbf{\color{cbScenario!50!black}provisional}
    \hfill {\color{cbInk!60}\footnotesize moderate $\cdot$ 55\%}\par\vspace{5pt}
    {\footnotesize\color{cbInk!70}\itshape What was President Gerald Ford's middle name?\par}\vspace{4pt}
    {\large Rudolph.\par}\vspace{3pt}
    {\footnotesize\color{cbInk!60} I'll reconsider this at my next learning cycle. Please ask again later for an updated answer.\par}};
  % DISCARD / conflicting
  \node[card, draw=cbNumbers!75!black, fill=cbNumbers!6] (c3) at (0,-3.4) {
    {\large\color{cbNumbers!70!black}$!$}\;\textbf{\color{cbNumbers!65!black}conflicting}
    \hfill {\color{cbInk!60}\footnotesize low $\cdot$ 40\%}\par\vspace{6pt}
    {\large\color{cbInk!45} (answer withheld)\par}\vspace{4pt}
    {\footnotesize\color{cbInk!60} Evidence exists but some signals disagree, independent verification is recommended.\par}};
  % ABSTAIN / insufficient
  \node[card, draw=cbPitfall!70!black, fill=cbPitfall!6] (c4) at (7.4,-3.4) {
    {\large\color{cbPitfall!65!black}$\times$}\;\textbf{\color{cbPitfall!60!black}insufficient}
    \hfill {\color{cbInk!60}\footnotesize very low}\par\vspace{6pt}
    {\large I do not know.\par}\vspace{4pt}
    {\footnotesize\color{cbInk!60} no supporting evidence was found.\par}};
\end{tikzpicture}
\caption[What the user actually sees]{The four verdicts rendered as user-facing cards. The
\emph{verified} and \emph{provisional} cards reproduce real deployed cycle-three records;
the \emph{conflicting} and \emph{insufficient} cards are representative of their verdict. A
\emph{verified} answer carries its
confidence; a \emph{provisional} answer is shown with a promise to re-check it; a
\emph{conflicting} verdict withholds the answer and explains that the evidence disagreed;
an \emph{insufficient} verdict returns an honest refusal. The answer and its reliability
signal always arrive together.}
\label{fig:cards}
\end{figure}

\section{Conflicting is not the same as insufficient}
\label{sec:output-failure-modes}

The most important subtlety in the envelope is that the two low-confidence verdicts,
discard and abstain, are \emph{different failure modes}, not two points on one
confidence scale. This was decided back at the gate (\Cref{ch:composite}) and it
surfaces here.

\begin{precise}[The split is by grounding, not by score]
Recall the gate: among answers below the deferral cut-point, the split between abstain
and discard is made by the grounding signal $\pgmax$, not by $\ustored$. An
\textbf{abstain} (tag \emph{insufficient}) means there was essentially no supporting
evidence at all, so the honest move is to say nothing. A \textbf{discard} (tag
\emph{conflicting}) means evidence existed but the signals disagreed about it, so the
answer is withheld with a different explanation. A consequence worth stating plainly:
a discarded answer can carry a \emph{higher} $\ustored$ than an abstained one. In the
deployed cycle-three run this is not a corner case but the norm: the discard and abstain
calibrated-confidence distributions overlapped almost entirely, with the chance that a
randomly chosen discarded answer scored higher than a randomly chosen abstained one close
to one half (near $0.46$). The two verdicts therefore genuinely cannot be ordered on a
single confidence axis; the split is by grounding, not by score. The two
tags are not ``a bit unsure'' versus ``very unsure''; they are ``the evidence
conflicts'' versus ``there is no evidence.'' Naming the failure mode, rather than only
its severity, is what lets the user respond appropriately, by rephrasing in one case
and by trusting the refusal in the other.
\end{precise}

\section{Deferral speaks in the system's own clock}
\label{sec:output-defer}

The deferred verdict has a caveat unlike the others, because deferral is a promise
about the future. A provisional answer is shown with a note that the system will
reconsider it, and the phrasing reflects how \caem{} actually updates: its cycle
boundary is triggered by \emph{query volume}, not by the wall clock, so the default
caveat is the cadence-neutral ``I'll reconsider this at my next learning cycle.'' A
deployment that updates on a known schedule can override this with a concrete cadence,
nightly or hourly or continuously, to give the user a real expectation. The deferral
message is the only place the user-facing surface acknowledges that the system learns
over time, which is the subject of the next chapter.

\section{Where the calibration mismatch is repaired}
\label{sec:output-repair}

\Cref{ch:problem} opened with the core hazard: a model states a falsehood with the
same confidence it states a truth, and a reader has no way to tell them apart. Stage 7
is the surface where that hazard is closed. Every served answer arrives with a
calibrated probability and a named reliability tag; a low-confidence answer is either
hedged, withheld, or replaced by an honest refusal; and the confidence number means
the same thing on every benchmark because of the calibration of \Cref{ch:composite}.
The confident lie no longer wears the same clothes as the confident truth, because the
clothes now include an honest label of how much to trust what was said.

\begin{defence}[If an examiner asks: ``why four outcomes rather than answer-or-refuse?'']
Because two of the four outcomes carry information a binary surface would discard.
\emph{Defer} separates the genuinely uncertain answer, which a later cycle's better
model may resolve, from the confidently wrong one; collapsing both into ``refuse''
would throw away a recoverable answer. And the \emph{conflicting}-versus-%
\emph{insufficient} split separates ``the evidence disagrees'' from ``there is no
evidence,'' two situations that call for different user responses. The four-outcome
contract is what lets the system distinguish \emph{wrong}, \emph{unknown}, and
\emph{not yet decided}, rather than flattening all three into a single refusal.
\end{defence}

\begin{bythenumbers}[The output contract]
\begin{itemize}[leftmargin=1.3em, itemsep=1pt]
  \item Two answer fields: raw (verbatim, for scoring) and display (curated, for the user).
  \item Display map: store/defer $\to$ stripped answer; abstain $\to$ ``I do not know.''; discard $\to$ withheld.
  \item Scaffold stripping shares the exact-match scorer's extraction routine.
  \item Envelope: calibrated percentage, semantic tag (verified / provisional / conflicting / insufficient), ordinal label, icon, show-or-hide flag, caveat.
  \item Discard versus abstain is split by grounding ($\pgmax$), not by $\ustored$.
\end{itemize}
\end{bythenumbers}

The per-query pipeline is now complete: a query has been triaged, routed, answered,
verified, scored, and presented with its confidence. Everything so far has happened
without the model learning anything. \Cref{ch:cycle} opens Part III and the second
clock, Stage 8, the cycle-boundary self-improvement loop, where the verified answers
this pipeline produced are folded back into the model and the whole system gets better.
